\cotefile{83}
\scanpage{96}
Supposons maintenant que
\[
\delta=\delta'=0\qquad \text{i.e.}\qquad b^2-4c=b'^2-4c'=0
\]
\unsure{i.e.} ce qui signifie \deleted{que \unsure{si on considère} $2$ alg. d'Az.} \inserted{que $A_1,A_2$ sont} loc. isom. à $k[T]/T^2$

% slanted
\margin{\uuline{faux} si on ne suppose pas que \unsure{l'on} est sur un corps, \uuline{ou} \unsure{$2$ inv.}}

(je l'ai peut-être vérifié dans le temps -- ça n'a pas l'air \deleted{évident} : fait évident, sur une base générale -- évident sur un corps, plus gén. si $\bigl[$On est ramené à ceci : si $b\in k$ tel que $b^2=0$, $\exists$ loc. fppf $\alpha,\beta\in k$, $\beta\in k^*$, tels que $\alpha^2=0$, $\beta(2\alpha+b)=0$ Evident si $2$ inv., mais faux si $2.1_k=0$ (car. $2$) car \deleted{\unsure{$\beta b=$|$\beta b=0$}} alors l'équation en $\beta$ donne $\beta b=0$ donc (comme $\beta\in k^{\unsure{*}}$) $b=0$, mais des situations $b^2=0$, $b\ne0$ ça \unsure{court|courre|croule|coince} ! $\bigr]$

\deleted{Ça signifie ainsi que} \deleted{Mais des alg. d'Az} \deleted{$A_1,A_2$ sont des sous-algèbres \unsure{correspondant|conjuguées}} \deleted{\inserted{normalis.}} \deleted{à la partie unipotente} \deleted{tirées d'un $\Theta\in M$ (Azumaya ambiant) qui vér.} \deleted{de Trace}
\[
\deleted{\operatorname{Tr}\Theta=\det\Theta=0}
\]
\deleted{Les \unsure{voici les normalisateurs|les unipotents} d'un tel} Il y a corr. $1-1$ entre une
% archive sheet number: 95
